\documentclass{article}
\usepackage{ifpdf} 
\usepackage{mla}
\title{Midterm Essay}
\author{Mackenzie Norman}
\begin{document}
\begin{mla}{Mackenzie}{Norman}{Lincoln}{ENG325}{\today}{essay}
\section*{Intro/Abstract}
\textit{Nyasha, Tambu and Lucia all are young women trapped within a colonial, racist, and patriarchal society. They are each affected by it differently, and in turn affect it differently}


\section*{The colonial weight}


\textit{``As if it is ever easy, And these days it is worse, with the poverty of blackness on one side and the weight of womanhood on the other''} Nervous Conditions; as it concludes, is a story of women. Women who are divided, complex, and despite all living in the same space, women who all affect it differently. Dangaremba says \textit{``People make individual choices. I think mapping the ground helps in making the choices. Such maps, written in an engaging way, are part of what I perceive some of my responsibility as a novelist to be.''} In viewing this novel as a map - or rather an atlas - the characters each as a map. The worldly Nyasha: a rebellious and neurotic youth, archetypal of the bildungsroman's hero and the main character Tambu: an obedient and poor child, with a deep yearning to be free of her colonial shackles.  Finally, Lucia: older and mysterious, with a seemingly magical ability to shape the world to her will. Continuing the metaphor the roads, the vessels that direct traffic through the map can be understood to be two dichotomies: food/hunger and education/naivete. This essay will inspect each characters path through these dichotomous roads, tracing their path, paying special importance to Lucia and her relative `success'


\section*{Nyasha, Oh Starving!}
\textit{``They do not like my language , my English, because it is authentic and my Shona, because it is not!'' }
\subsection*{Nyasha's interaction with colonialism}
Nyasha is deeply aware of how colonialism affects her life. She uses food as a metaphor since she has "tasted" what it feels like to not be in Africa. \textit{``Nyasha tucked into the vegetables with the rest of us. When my mother offered her the sour milk she had asked for, she became very morose. She refused to eat'' } She is the only character in the book who is  \textbf{outwardly} Anticolonial
\subsection*{ Her Relation to food}
Nyasha has anorexia. Food to Nyasha is what her father uses to control her. An unwanted social lubricant? As above, food is often a reminder of her place in colonial africa, the food her parents serve and eat is a facsimile of what she ate in England, akin to her in authentic life back in africa. While african food - like her shona is not authentic to her.
\subsection*{ Her relation to education}
Nyasha is one of the more educated people and clearly one of the most discerning. She often verges on too smart for her own good. 


\section*{ Tambu forces food down}
\subsection*{ Tambu's interaction with colonialism}
Tambu has an odd relationship with colonialism in that she is obviously one of the most affected by it in the book, 
\subsection*{ Her Relation to food}
Tambus relation to food mirrors her relation to education. Food is a means to escape. We see this literally in the beginning when she is able to pay for her schooling by selling food. She differs from Nyasha in that african food is authentic to her. Compared to Babamakurus whole family even she was the only one who liked sadza \textit{``I was very pleased to see the sadza when it came, although nobody else seemed to care for it. This was embarrassing. There were many things that were embarrassing about that meal: my place looked as though a small and angry child had been fed there; here I was with a spoon in my hand instead of a fork and now Maiguru was dishing sadza on to my plate, sadza that nobody else would eat'' }


\subsection*{ Her relation to education}
Tambu loves learning, for `selfish' reasons (nearly everyone in this book is `selfish') Tambu believes that if she is educated enough she will not have to marry, or at least will have far more freedom than she would if she was not educated. She enjoys education and takes to it well. Connecting to food, she does not turn it down, and instead uses it to hopefully get her where she wants to be.

\section*{ Lucia, a middle ground?}

\subsection*{ Her Relation to food}
Lucia has so many oddly masculine attributes. She is able to provide not just food, but meat.  and is somehow magically always well fed. 
\subsection*{ Her Relation to education}
\end{mla}
\end{document}